% =============================================================================
\section{Méthodes d'optimisation retenues}
\label{sec:methods}
% =============================================================================

Cette section présente les trois familles de métaheuristiques retenues après
une phase exploratoire (cf.\ section~\ref{sec:autres_algorithmes}).  Pour
chacune, nous rappelons le principe, donnons les équations clés et précisons
les hyperparamètres.  Le choix de ces trois familles est motivé par leur
complémentarité : CMA-ES excelle sur les problèmes corrélés (prob1--prob3),
PSO sur les paysages fortement multimodaux (prob4), et DE/SHADE offre un
compromis robuste avec auto-adaptation des paramètres.

% ----- DE --------------------------------------------------------------------
\subsection{Évolution différentielle (DE)}
\label{sec:de}

\subsubsection{Principe}

L'évolution différentielle~\cite{storn1997} exploite les \emph{différences}
entre individus de la population comme vecteur de perturbation.  La variante
retenue est \texttt{DE/current-to-best/1/bin}.

\subsubsection{Mutation}

Pour chaque individu cible $\mathbf{x}_i$, un vecteur mutant est construit :
\begin{equation}
\label{eq:de-mutation}
  \mathbf{v}_i = \mathbf{x}_i + F_i \cdot (\mathbf{x}_{\text{best}} -
  \mathbf{x}_i) + F_i \cdot (\mathbf{x}_{r_1} - \mathbf{x}_{r_2}),
\end{equation}
où $r_1 \neq r_2 \neq i$ sont tirés aléatoirement et $\mathbf{x}_{\text{best}}$
est le meilleur individu courant.

\subsubsection{Croisement binomial}

Le vecteur d'essai $\mathbf{u}_i$ est construit par :
\[
  u_{i,j} =
  \begin{cases}
    v_{i,j}, & \text{si } \text{rand}_j < CR_i \text{ ou } j = j_{\text{rand}}, \\
    x_{i,j}, & \text{sinon},
  \end{cases}
\]
où $j_{\text{rand}}$ est un indice choisi aléatoirement pour garantir qu'au
moins une composante provient du mutant.

\subsubsection{Sélection gloutonne}

L'individu d'essai remplace la cible si et seulement si sa \textit{fitness} est
meilleure ou égale :
$\mathbf{x}_i^{(g+1)} = \mathbf{u}_i$ si $f(\mathbf{u}_i) \leq
f(\mathbf{x}_i^{(g)})$, sinon $\mathbf{x}_i^{(g+1)} = \mathbf{x}_i^{(g)}$.

\subsubsection{Auto-adaptation jDE}

Les paramètres $F_i$ et $CR_i$ sont adaptés individuellement à chaque
génération~\cite{brest2006} :
\[
  F_i^{(g+1)} =
  \begin{cases}
    F_l + \text{rand} \cdot F_u, & \text{avec probabilité } \tau_1 = 0{,}1, \\
    F_i^{(g)}, & \text{sinon},
  \end{cases}
  \qquad
  CR_i^{(g+1)} =
  \begin{cases}
    \text{rand}, & \text{avec probabilité } \tau_2 = 0{,}1, \\
    CR_i^{(g)}, & \text{sinon},
  \end{cases}
\]
avec $F_l = 0{,}1$ et $F_u = 0{,}9$.  Ce mécanisme supprime le besoin de
régler manuellement $F$ et $CR$.

\subsubsection{Recherche locale hybride}

Toutes les 5~générations, un (1+1)-ES gaussien est appliqué au meilleur
individu pendant 100~évaluations au maximum, avec un pas initial de 0,5 et
la règle d'adaptation dite « du 1/5 de succès ».

\paragraph{Paramètres retenus.}
$N_P = 10d$, $F_{\text{init}} = 0{,}7$, $CR_{\text{init}} = 0{,}9$,
$\tau_1 = \tau_2 = 0{,}1$.

% ----- CMA-ES ----------------------------------------------------------------
\subsection{CMA-ES (\emph{Covariance Matrix Adaptation — Evolution Strategy})}
\label{sec:cmaes}

\subsubsection{Principe}

CMA-ES~\cite{hansen2001,hansen2016} est une stratégie évolutionnaire qui
échantillonne $\lambda$ descendants à partir d'une distribution multinormale
$\mathcal{N}(\mathbf{m}, \sigma^2 \mathbf{C})$ et met à jour ses paramètres
--- le centroïde $\mathbf{m}$, le pas global $\sigma$ et la matrice de
covariance $\mathbf{C}$ --- à partir des $\mu$ meilleurs individus.

\subsubsection{Échantillonnage}

À la génération $g$, les $\lambda$ descendants sont générés par :
\begin{equation}
\label{eq:cma-sample}
  \mathbf{x}_k^{(g+1)} = \mathbf{m}^{(g)} + \sigma^{(g)} \cdot
  \mathbf{B}\,\text{diag}(\mathbf{d})\,\mathbf{z}_k,
  \qquad \mathbf{z}_k \sim \mathcal{N}(\mathbf{0}, \mathbf{I}),
\end{equation}
où $\mathbf{C} = \mathbf{B}\,\text{diag}(\mathbf{d}^2)\,\mathbf{B}^T$ est la
décomposition en valeurs propres de la matrice de covariance.

\subsubsection{Mise à jour du centroïde}

Les $\mu$ meilleurs descendants (triés par \textit{fitness} croissante) sont combinés
par moyenne pondérée :
\begin{equation}
  \mathbf{m}^{(g+1)} = \sum_{i=1}^{\mu} w_i \, \mathbf{x}_{i:\lambda}^{(g+1)},
  \qquad \sum_{i=1}^{\mu} w_i = 1, \quad w_1 \geq \ldots \geq w_\mu > 0.
\end{equation}

\subsubsection{Chemins d'évolution cumulatifs}

Deux chemins d'évolution sont maintenus pour contrôler le pas~$\sigma$ et la
matrice~$\mathbf{C}$ :
\begin{align}
  \mathbf{p}_\sigma^{(g+1)} &= (1 - c_\sigma)\,\mathbf{p}_\sigma^{(g)}
    + \sqrt{c_\sigma(2 - c_\sigma)\,\mu_{\text{eff}}}\;
      \mathbf{C}^{-1/2}\,\frac{\mathbf{m}^{(g+1)} - \mathbf{m}^{(g)}}{\sigma^{(g)}},
    \label{eq:ps} \\
  \mathbf{p}_c^{(g+1)} &= (1 - c_c)\,\mathbf{p}_c^{(g)}
    + h_\sigma\,\sqrt{c_c(2 - c_c)\,\mu_{\text{eff}}}\;
      \frac{\mathbf{m}^{(g+1)} - \mathbf{m}^{(g)}}{\sigma^{(g)}},
    \label{eq:pc}
\end{align}
avec $\mu_{\text{eff}} = 1/\sum_{i=1}^{\mu} w_i^2$ la variance effective des
poids.  L'indicateur $h_\sigma$ vaut 1 si
$\|\mathbf{p}_\sigma^{(g+1)}\| < \bigl(1{,}4 + 2/(d+1)\bigr)\,\chi_d$ et 0
sinon ; il évite une mise à jour de rang~1 lorsque $\sigma$ est trop petit
(stagnation de $\mathbf{p}_\sigma$).

\subsubsection{Adaptation du pas $\sigma$ (CSA)}

Le \textit{Cumulative Step-size Adaptation} (CSA) ajuste $\sigma$ en comparant
la norme de $\mathbf{p}_\sigma$ à sa valeur attendue sous marche aléatoire
isotrope :
\begin{equation}
\label{eq:sigma-update}
  \sigma^{(g+1)} = \sigma^{(g)} \cdot
  \exp\!\left(\frac{c_\sigma}{d_\sigma}
    \left(\frac{\|\mathbf{p}_\sigma^{(g+1)}\|}{\chi_d} - 1\right)\right),
\end{equation}
où $d_\sigma = 1 + 2\max\!\bigl(0,\,\sqrt{(\mu_{\text{eff}}-1)/(d+1)} -
1\bigr) + c_\sigma$ est le facteur d'amortissement~\cite{hansen2001} et
$\chi_d \approx \sqrt{d}\,(1 - 1/(4d) + 1/(21 d^2))$ l'espérance de
$\|\mathcal{N}(\mathbf{0}, \mathbf{I})\|$.

\emph{Discussion.}  Si les pas successifs sont alignés (bonne direction), la
norme de $\mathbf{p}_\sigma$ croît et $\sigma$ augmente, accélérant la
convergence.  À l'inverse, des pas incohérents font décroître $\sigma$, ce qui
ralentit l'exploration --- un mécanisme d'auto-adaptation fondamental.

\subsubsection{Adaptation de la matrice de covariance $\mathbf{C}$}

La covariance est mise à jour par une combinaison de rang 1 et rang $\mu$ :
\begin{equation}
\label{eq:C-update}
  \mathbf{C}^{(g+1)} = (1 - c_1 - c_\mu)\,\mathbf{C}^{(g)}
    + c_1\,\mathbf{p}_c \mathbf{p}_c^T
    + c_\mu \sum_{i=1}^{\mu} w_i \,
      \frac{(\mathbf{x}_{i:\lambda} - \mathbf{m}^{(g)})
            (\mathbf{x}_{i:\lambda} - \mathbf{m}^{(g)})^T}
           {\sigma^{(g)2}}.
\end{equation}

\emph{Discussion.}  C'est cette adaptation de $\mathbf{C}$ qui distingue
fondamentalement CMA-ES de GA et DE : l'algorithme apprend la \emph{structure de corrélation} du paysage de \textit{fitness}, ce qui lui permet de concentrer la
    recherche dans les directions les plus prometteuses.

% ----- PSO -------------------------------------------------------------------
\subsection{Optimisation par essaim particulaire (PSO)}
\label{sec:pso}

\subsubsection{Principe}
L'optimisation par essaim particulaire~\cite{kennedy1995} s'inspire du
comportement collectif des bancs de poissons ou des volées d'oiseaux.  Chaque
particule $i$ se déplace dans l'espace de recherche $\mathbb{R}^d$ avec une
vitesse $\mathbf{v}_i$ mise à jour en fonction de sa meilleure position
personnelle $\mathbf{p}_i$ (\textit{personal best}) et de la meilleure position
globale de l'essaim $\mathbf{g}$ (\textit{global best}).

\subsubsection{Mise à jour de la vitesse et de la position}
À chaque itération $t$, la vitesse et la position sont mises à jour selon :
\begin{align}
\label{eq:pso-velocity}
    \mathbf{v}_i^{(t+1)} &= w \, \mathbf{v}_i^{(t)}
      + c_1 \, \mathbf{r}_1 \odot (\mathbf{p}_i - \mathbf{x}_i^{(t)})
      + c_2 \, \mathbf{r}_2 \odot (\mathbf{g} - \mathbf{x}_i^{(t)}), \\
\label{eq:pso-position}
    \mathbf{x}_i^{(t+1)} &= \mathbf{x}_i^{(t)} + \mathbf{v}_i^{(t+1)},
\end{align}
où $w$ est le coefficient d'inertie, $c_1$ et $c_2$ les coefficients
d'accélération (cognitif et social), $\mathbf{r}_1, \mathbf{r}_2 \sim
\mathcal{U}(0,1)^d$ des vecteurs aléatoires composante par composante, et
$\odot$ le produit de Hadamard.

\subsubsection{Gestion de l'inertie et mécanisme anti-stagnation}
Pour favoriser l'exploration en début de recherche et l'exploitation en fin,
l'inertie décroît linéairement de $w_{\max} = 0{,}9$ à $w_{\min} = 0{,}4$
au cours des itérations.  Un mécanisme d'anti-stagnation complète cette
dynamique : si $\mathbf{g}$ ne s'améliore pas pendant $10N$ générations
consécutives, 50\,\% des particules sont réinitialisées aléatoirement afin
de relancer la diversité de l'essaim.

\emph{Discussion.}  PSO se distingue des stratégies d'évolution (CMA-ES, DE)
par l'absence de modèle probabiliste explicite de la distribution des
solutions.  La dynamique vitesse-position confère à l'essaim une mémoire
cinétique qui lui permet de traverser des bassins d'attraction étroits,
propriété particulièrement avantageuse sur les paysages labyrinthiques
(cf.\ section~\ref{sec:prob4}).

\paragraph{Paramètres retenus.}
Taille de l'essaim $N = 10d$, $c_1 = c_2 = 2{,}0$, $v_{\max} = 0{,}2
\times \text{range}$, seuil de stagnation $10N$ générations.
